\section*{الفصل \ref{ch:b3:ocean-dynamics} --- ديناميكا المحيطات}

\begin{solution}{exo:b3:ocean-dynamics:1}
$f = \qty{9.37e-5}{s^{-1}}$; $\tau/(\rho f) = 0.12/(1025\times\num{9.37e-5})
= \qty{1.25}{m^2/s}$، على يمين الريح: جنوبًا. عبر
\qty{1000}{km}، \qty{1.25e6}{m^3/s}، \qty{1.25}{\sverdrup}.
\end{solution}

\begin{solution}{exo:b3:ocean-dynamics:2}
$f = \qty{1.031e-4}{s^{-1}}$; $D_E = \pi\sqrt{2\times0.05/\num{1.031e-4}}
= \qty{98}{m}$. يساوي التيار هناك $\mathrm e^{-\pi} = \num{0.043}$ من
تيار السطح، ويجري في الاتجاه المعاكس.
\end{solution}

\begin{solution}{exo:b3:ocean-dynamics:3}
$\beta = 2\Omega\cos\ang{30}/a = \qty{1.98e-11}{m^{-1}.s^{-1}}$; $V =
\num{-1e-7}/(1025\times\num{1.98e-11}) = \qty{-4.92}{m^2/s}$؛ عبر
\qty{5000}{km}، \qty{-24.6}{\sverdrup}: جنوبًا.
\end{solution}

\begin{solution}{exo:b3:ocean-dynamics:4}
يعود المكوّنان إلى الطور نفسه بعد $1/(1/12 - 1/12.4206) =
\qty{354}{h}$، \num{14.77}\text{ يومًا}، نصف الشهر الاقتراني.
\end{solution}

\begin{solution}{exo:b3:ocean-dynamics:5}
عند \ang{40} شمالًا، $f = \qty{9.37e-5}{s^{-1}}$، $\beta =
\qty{1.75e-11}{m^{-1}.s^{-1}}$، $\tau_x = \qty{0.05}{Pa}$
و$\partial\tau_x/\partial y = 0.1/\num{1.11e6} = \qty{9.0e-8}{Pa/m}$. الحدّ
الأول، $-\num{9.0e-8}/(1025\times\num{9.37e-5}) =
\qty{-9.38e-7}{m/s}$؛ الحدّ الثاني، $0.05\times\num{1.75e-11}/(1025\times
(\num{9.37e-5})^2) = \qty{+9.7e-8}{m/s}$. $w_E = \qty{-8.4e-7}{m/s}$،
\qty{-26.5}{m/yr}: يُدفع الماء إلى الأسفل.
\end{solution}

\begin{solution}{exo:b3:ocean-dynamics:6}
$\delta_S = \num{2e-6}/\num{2e-11} = \qty{100}{km}$؛ السرعة $=
\num{30e6}/(\num{1e5}\times800) = \qty{0.375}{m/s}$.
\end{solution}

\begin{solution}{exo:b3:ocean-dynamics:7}
$\beta = \qty{1.98e-11}{m^{-1}.s^{-1}}$; $\delta_M = (10^4/\num{1.98e-11})^{1/3}
= \qty{79.6}{km}$، والتيار النفاث عرضه $3.628\,\delta_M = \qty{289}{km}$.
مع $A_H$ أصغر ثماني مرّات ينخفض كلٌّ منهما إلى النصف، لأنّ $8^{1/3} = 2$.
\end{solution}

\begin{solution}{exo:b3:ocean-dynamics:8}
$f = \qty{8.57e-5}{s^{-1}}$; $v = (21.3 - 14.5)/(\num{8.57e-5}\times\num{8e4})
= \qty{0.99}{m/s}$، موجبة مع $B$ على اليمين: بالنظر من $A$ إلى $B$،
شرقًا، يتّجه التيار شمالًا. يعلو سطح البحر
$6.8/9.81 = \qty{0.69}{m}$ أعلى عند $B$؛ عمود الماء عند $B$ هو
الأخفّ والأدفأ، على يمين التيار.
\end{solution}

\begin{solution}{exo:b3:ocean-dynamics:9}
مدى المدّ الربيعي $1.5K(1 + 0.46) = \qty{0.78}{m}$؛ مدى المدّ النيابي $1.5K(1 - 0.46)
= \qty{0.29}{m}$؛ النسبة $1.46/0.54 = 2.7$.
\end{solution}

\begin{solution}{exo:b3:ocean-dynamics:10}
كما في \cref{thm:b3:ocean-dynamics:kelvin}، $u = \pm(g/c)\eta$ من أجل $F(x \mp
ct)$، و$fu = -g\,\partial\eta/\partial y$ يصبح $\beta_0y(\pm g/c)G =
-gG'$، إذن $G = \exp(\mp\beta_0y^2/2c)$. لا يُحصر على جانبي خط الاستواء
إلا الإشارة العليا، أي الموجة المتجهة شرقًا. $\beta_0 =
\qty{2.29e-11}{m^{-1}.s^{-1}}$، $\sqrt{2.5/\num{2.29e-11}} = \qty{330}{km}$.
\end{solution}

\begin{solution}{exo:b3:ocean-dynamics:11}
(أ)~بتكامل $\beta v = f\,\partial w/\partial z$ من القاع، $w = 0$،
حتى أعلى الطبقة اللُّجّية، $w = W_0$: $\beta V = fW_0$، و$f/\beta =
a\tan\varphi$. (ب)~$V = \num{6.371e6}\times0.577\times10^{-7} =
\qty{0.368}{m^2/s}$، \qty{1.84}{\sverdrup} عبر \qty{5000}{km}، نحو القطب.
(ج)~من تيار الحدّ الغربي العميق الذي تغذّيه المنابع في خطوط العرض العالية؛
يجري الجزء الداخلي نحو المنابع، ويجري تيار الحدّ بعيدًا
عنها.
\end{solution}

\begin{solution}{exo:b3:ocean-dynamics:12}
(أ)~مع $z$ موجّهًا نحو الأعلى انطلاقًا من القاع، $A_zW'' = \mathrm if(W - W_g)$، إذ
يكون تدرّج الضغط هو تدرّج الجريان الجيوستروفي الذي فوقه. $W - W_g =
C\mathrm e^{-(1+\mathrm i)z/\delta}$ محدودة نحو الأعلى، و$W(0) = 0$ يعطي
$C = -W_g$. (ب)~عجز النقل هو $\int_0^\infty(W - W_g)\,\mathrm dz
= -W_g\delta/(1+\mathrm i) = -\tfrac12\delta(1 - \mathrm i)W_g$، أي
$-\tfrac12\delta(u_g + v_g,\ v_g - u_g)$. ولأنّ الجريان الجيوستروفي عديم
التباعد، $w = -\nabla\cdot(\text{العجز}) = \tfrac12\delta(\partial_x
v_g - \partial_y u_g) = \tfrac12\delta\zeta_g$. (ج)~$\delta =
\sqrt{0.02/10^{-4}} = \qty{14.1}{m}$; $r = 10^{-4}\times14.1/8000 =
\qty{1.77e-7}{s^{-1}}$، زمن تباطؤ الدوران \num{65}\text{ يومًا}.
\end{solution}

\begin{solution}{pb:b3:ocean-dynamics:1}
\textbf{1.} $\tau_x = \qty{-0.1}{Pa}$ (شرقية)، 0 و\qty{+0.1}{Pa}
(غربية).

\textbf{2.} $f = \num{4.99e-5}$، \num{8.37e-5}
و\qty{1.117e-4}{s^{-1}}.

\textbf{3.} $V_E = -\tau_x/(\rho f)$: \qty{+1.96}{m^2/s} (شمالًا) عند
\ang{20} شمالًا و\qty{-0.87}{m^2/s} (جنوبًا) عند \ang{50} شمالًا. ويتقاربان.

\textbf{4.} $D_E = \pi\sqrt{2\times0.05/\num{8.37e-5}} = \qty{109}{m}$.

\textbf{5.} $\tau_x = 0$ عند \ang{35} شمالًا، إذن $w_E = -(1/\rho f)\,
\partial\tau_x/\partial y = -\pi\tau_0/(\rho fL) = \qty{-1.10e-6}{m/s}$،
\qty{-34.7}{m/yr}، إلى الأسفل.

\textbf{6.} $500/34.7 = \num{14.4}\text{ سنة}$.

\textbf{7.} $L = a\times\pi/6 = \qty{3336}{km}$; $\beta =
\qty{1.875e-11}{m^{-1}.s^{-1}}$.

\textbf{8.} $\operatorname{curl}_z\vec\tau = -\pi\tau_0/L =
\qty{-9.42e-8}{Pa/m}$.

\textbf{9.} $V = \num{-9.42e-8}/(1025\times\num{1.875e-11}) =
\qty{-4.90}{m^2/s}$، جنوبًا.

\textbf{10.} $4.90\times\num{6e6} = \qty{2.94e7}{m^3/s}$،
\qty{29.4}{\sverdrup} جنوبًا.

\textbf{11.} $\qty{4.9}{mm/s}$؛ \qty{1000}{km} في \qty{2.0e8}{s}،
\num{6.5}~سنة.

\textbf{12.} $fw_E/\beta = \num{8.37e-5}\times(\num{-1.10e-6})/
\num{1.875e-11} = \qty{-4.90}{m^2/s}$: القيمة نفسها، لأنّ نقل إيكمان
ينعدم حيث $\tau_x = 0$. عمود مضغوط بالمعدّل $w_E$
يحتفظ بدوّاميته الكامنة $f/h$ فقط بالانتقال نحو $f$ أصغر، جنوبًا،
بالمعدّل $\beta V_g = fw_E$.

\textbf{13.} لا يمكن أن تكون الطبقة الحدّية إلّا طبقةً يتلاشى فيها التصحيح المطبَّق على
الداخل بالابتعاد عن الجدار؛ ومع $\beta > 0$ يكون ذلك
$\mathrm e^{-x/\delta_S}$، وهو يتلاشى شرقًا، انطلاقًا من الجدار الغربي
فقط. بلغة الدوّامية، يزيل الاحتكاك في تيار متّجه شمالًا ضد جدار
شرقيّ يزيل الدوّامية الموجبة، وهي الإشارة الخاطئة.

\textbf{14.} $\delta_S = \num{e-6}/\num{1.875e-11} = \qty{53.3}{km}$.

\textbf{15.} $A = \qty{4.90}{m^2/s}$، $X/\delta_S = 112.5$:
$4.90\times111.5 = \qty{546}{m^2/s}$، \qty{0.55}{m/s} عبر
\qty{1000}{m}.

\textbf{16.} حيث $(X/\delta_S)\mathrm e^{-x/\delta_S} = 1$: $x =
\delta_S\ln(X/\delta_S) = \qty{252}{km}$.

\textbf{17.} $A(X - \delta_S - x_0) = 4.90\times(6000 - 53 - 252)\times
10^3 = \qty{27.9}{\sverdrup}$. يستمرّ التدفّق الداخلي جنوبًا في
الشريط المجاور للجدار؛ أمّا النقص، $A(x_0 + \delta_S) =
\qty{1.5}{\sverdrup}$، فهو النقل الداخلي الذي كان سيعبر
ذلك الشريط.

\textbf{18.} $\delta_M = (\num{5e3}/\num{1.875e-11})^{1/3} = \qty{64.4}{km}$;
التيار النفّاث $3.628\times64.4 = \qty{233}{km}$.

\textbf{19.} $1.163\times29.4 - A\times\qty{233}{km} = 34.2 - 1.1 =
\qty{33.0}{\sverdrup}$؛ الفائض، \qty{16}{\percent} من النقل الداخلي
يعود جنوبًا في التيار المعاكس مباشرةً شرق التيار النفّاث.

\textbf{20.} $c = \sqrt{9.81\times4000} = \qty{198}{m/s}$; $L_R =
198/\num{8.37e-5} = \qty{2370}{km}$.

\textbf{21.} $198\times12.4206\times3600 = \qty{8860}{km}$.

\textbf{22.} محيط الحوض $2(6000 + 3336) = \qty{18670}{km}$، \qty{26.2}{h}،
\mbox{\textdir TLT 2.1 M$_2$} دورات.

\textbf{23.} الدائرة هي $2\pi a\cos\ang{35} = \qty{32790}{km}$، تُقطع
في يوم قمري، \qty{24.84}{h}: \qty{367}{m/s}، وهذا يتطلّب $c^2/g =
\qty{13.7}{km}$ من الماء.

\textbf{24.} مع نصف قطر روسبي مماثل لعرض الحوض،
وطول موجة \mbox{\textdir TLT M$_2$} مماثلًا لطوله، تقريبًا نقطة أمفيدرومية واحدة
قرب مركز الحوض؛ يدور المدّ المرتفع عكس اتّجاه عقارب الساعة
حولها، مُبقيًا الساحل على يمينه.

\textbf{25.} $\tfrac34\times0.357\times\cos^2\ang{35} = \qty{0.18}{m}$.

\textbf{26.} يُعيد تيار الحدّ الغربي شمالًا
\textbf{\qty[text-series-to-math=true]{29}{\sverdrup}} الذي يحمله الداخل
جنوبًا (\qty{29.4}{\sverdrup} وفق الحساب).
\end{solution}
